% !TeX root = ../navier-stokes-manuscript.tex
% ============================================================
% 1.3 Gold Standard Proof
% ============================================================
\subsection{Gold Standard Proof}\label{sub:GoldStandardProof}

The direct proof is short once one estimate is available.
For some fixed \(s>5/2\), derive from \(u_0\), \(\nu\), and the Navier--Stokes equations that, for every finite \(T>0\),
\[
  \sup_{0\leq t<\min\{T,T_*\}}
  \lVert u(t)\rVert_{H^s(\mathbb R^3)}
  \leq
  C_s(u_0,\nu,T)
  <
  \infty.
\]
This is the missing Gold-standard estimate.

Its consequence follows from the continuation argument above.
Sobolev embedding controls \(\nabla u\) in \(L^\infty\).
The differentiated equations then propagate every fixed higher \(H^m\) norm from the smooth datum, the pressure equation recovers the spatial derivatives of \(p\), and the momentum equation recovers the time derivatives.
If \(T_*<\infty\), take \(T=T_*\) in the estimate and restart local theory from a sufficiently late time slice.
The uniform \(H^s\) bound gives enough lifespan to cross \(T_*\), contradicting the definition of maximal time.
This bound forces
\[
  T_*=\infty.
\]

The difficulty lies entirely in deriving that bound from the controls already supplied by the equations.
Energy gives
\[
  \int_0^T
  \lVert\nabla u(t)\rVert_{L^2}^{2}\,dt
  <
  \infty.
\]
At the first differentiated level, the nonlinear term leads after interpolation and viscous absorption to
\[
  \frac{d}{dt}
  \lVert\nabla u\rVert_{L^2}^{2}
  +
  \nu\lVert\Delta u\rVert_{L^2}^{2}
  \leq
  C\nu^{-3}
  \lVert\nabla u\rVert_{L^2}^{6}.
\]
The available time integral controls the second power of the gradient.
The differentiated inequality asks for control strong enough to handle its sixth power.
Short high peaks can satisfy the first statement and defeat the second, so this estimate stops before the first derivative level closes.
That observation tests the strength of the estimate; it does not claim that the fluid produces such peaks.

A Gold-standard proof has to obtain the continuation-grade bound from transport, pressure, incompressibility, and viscosity acting in the original equations.
Once that estimate is proved, the rest of the argument is already in place.
